\section{Conclusions}\label{sec:conclusions}

We have decomposed the bimetric benchmark of \citet{Hogas2025} probe by probe, and profiled the anchor to separate properties of the theory from properties of the chosen point.

\begin{enumerate}
\item \textbf{BAO.} At the benchmark and across the reported fixed-anchor grids, DESI~DR2 BAO differences are small: $\Dchi=0.179$ with the full covariance. ``Blindness'' is limited to $(\alpha,\beta,\mFP,\theta)$ in this fixed-anchor statistical domain and does not extend to $\OL$ or across the anchor plane.
\item \textbf{Mechanism.} At the adopted anchor the late-time background is \LCDM{} with the effective density $\Om^{\rm eff}=\cos^2\!\theta\,\Om$ to a maximum relative residual of $3.4\times10^{-5}$ in $E^2$ over $z\le2.33$ (verified), the dark degeneracy. Fixed-anchor BAO blindness and SNe density sensitivity are the two faces of that single fact.
\item \textbf{SNe.} Pantheon+ ($z_{\rm HD}$ convention) gives $\Dchi=2.94$ with the full STAT+SYS covariance and $7.32$ with the diagonal covariance. These fixed-benchmark, per-probe differences do not constitute a joint model-selection significance.
\item \textbf{The anchor.} Under the adopted $\lambda=0$ CMB protocol, the descriptive sum of per-probe $\Dchi$ values is flat in $\theta$ from $0$ to $6\deg$. Its scanned minimum is $2.98$ at $(5\deg,0.6951)$. The profiled benchmark-angle point ($6.1\deg$, $0.6976$) gives $3.04$, or $0.06$ above the minimum; the adopted benchmark gives $3.115$, or $0.14$ above it. The GR-side floor is $3.91$, about $0.9$ above the minimum, showing residual inter-probe tension rather than a unique failure of quartic closure. At the adopted benchmark $2.94$ of the $3.115$ floor belongs to SNe. This sum is not a joint likelihood; its robust content is the map of competing probe pulls and parameter mismatch. The strongest discriminating diagnostic among those considered is the scalar sector (item~7).
\item \textbf{$\Om$ rigidity.} The near-rigidity of $\Om$ under $(\alpha,\beta)$ holds across the 19-point grid; a global impossibility claim over the full interaction space is not made.
\item \textbf{Compressed CMB.} Within this compressed diagnostic, the one-parameter transplant family $R(\lambda,\theta)$ shows structural saturation with $\lambda$ free and therefore no robust $\theta$ information beyond the $10^{-5}$--$3.6\times10^{-3}$ effective-law residual (\cref{sec:cmbres}); the five-point sweep and local parabolas do not certify precise continuous vertices. Under the adopted $\lambda=0$ protocol (true density in $R$, the $\lambda$-systematic quoted) the benchmark fits at $\chi^2=0.0001703$ and large angles are penalised; a principled re-derivation awaits a Boltzmann-level treatment, which partial perturbation calculations (e.g.\ tensor modes) do not substitute for.
\item \textbf{Scalar sector.} On the adopted finite branch and linear/WKB diagnostic the instability boundary lies at $z_*=7.8771$, and every tabulated point of the scanned anchor set is scalar-unstable at $\mFP/H_0=25$; the self-similar reduction explains the scanned crossings; exact uniqueness is proved only for the benchmark finite branch, not over the full parameter space. The formal growth exponent ($\log_{10}\Gamma=87.87$ at $k=0.05\ \mathrm{Mpc}^{-1}$, $h=0.700$) accumulates predominantly in the screened regime and is not a physical growth estimate; for the adopted order-of-magnitude $z_V=11.72$, the unscreened contribution alone gives $\log_{10}\Gamma_u=10.8$, subject to the same diagnostic restrictions. No scale simultaneously satisfies the illustrative $\Gamma\le10$ cut and the applicability requirements of the full-window diagnostic (\cref{tab:dilemma}), but this does not physically exclude low-$k$ modes. In the prescribed cosmology of \cref{tab:mass}, $\mFP/H_0\gtrsim4.1\times10^4$ gives $z_*>1090$ and hence no gradient-unstable interval at $z<1090$; the earlier interval $z>z_*$ remains unstable. Physical viability at other points or after nonlinear evolution remains undecided.
\item \textbf{Self-similarity.} The boundary is universal in $u=\Om(1+z)^3/(\sin^2\!\theta\,\mFP^2)$ at fixed $t$ ($w=tu$ across angles) and collapses into the single master equation $\Om(1+z_*)^3=\sin^2\!\theta\,\mFP^2\,u_*(\varepsilon)$: $(1+z_*)\propto\mFP^{2/3}$ (slope $0.6665$ against $2/3$), the $\Lambda$ correction, and $H^2(z_*)=c^2\mFP^2+\kappa\,\OL$ with branch constants $c=0.582199$, $\kappa=0.374329$ are one statement (\cref{app:selfsim}). An exact algebraic certificate for the reduced $\varepsilon=0$ benchmark branch, together with a full-model numerical $\dd y/\dd N$ scan, rules out a Higuchi ghost within their stated scopes; the identified pathology on this branch is the scalar gradient instability. This is not a full stability certification. The law is a within-branch statement of the self-similar regime ($\varepsilon\ll1$), not a theorem over the full parameter space.
\end{enumerate}

The combined result distinguishes fit quality from viability: the H\&M joint-fit preference is not reproduced here as a joint likelihood, our per-probe decomposition explains the background response, and the scalar diagnostic shows why a good background fit alone establishes neither viability of the point nor preference for the underlying theory.
